\subsection{Avancements}
Au cours de la journée, l'architecture embarquée a majoritairement été commencé pour retrouver les services principaux.
L'utilisation d'un seul \textit{Shared memory region} pour représenter le robot pour tout les services à été choisi pour des questions de simplicités.
Le système Ethernet permettait d'entrer sur le robot de manière fiable et efficace.
En utilisant un utilitaire \textit{SCP} pour envoyer les services à des fins de déverminage, l'exécution de \textit{scripts} permettant d'automatiser le deployement de l'environement de test permettait de rapidement avoir un prototype testable.
Il a été possible d'évaluer des constantes pour calibrer le capteur de ligne afin que le robot puisse normaliser et moyenner, par le biais d'algorithmes, une ligne.

\subsection{Modification}
Afin de démarrer le robot et le WIFI sans à avoir à utiliser le cable ethernet, la proposition décrit l'utilisation de deux boutons et l'intégration d'un service \textbf{PeripheralsService} pour détecter l'intéraction avec les boutons et exécuter l'action suivant lesdites intéractions.
De plus, un système sonore permetterais de confirmer que le robot est démarré ou bien de connaître l'état des systèmes (WIFI ouvert/fermé, démarrage de la séquence, etc.)
